\documentclass{ximera}
\begin{document}

    \title{Lecture 18}
    
    \begin{abstract} 
        - Review Gaussian elimination and elementary matrices \\
        - LU decomposition
    \end{abstract}
    \maketitle

\textit{This lecture is an alternative to the first 2/3 of Lecture 20 in Trefethen \& Bau.}

Recall that the process of Gaussian elimination, to put a matrix into REF or RREF, are given by the following 3 rules:
\begin{enumerate}
    \item  Interchange two equations.
    \item  Multiply an equation by a nonzero constant.
    \item  Add a constant multiple of one equation to another.
\end{enumerate}

Let's consider the system of equations 
\begin{align*}
    x+2y-z &= 3 \\
    2x+5y+z &=12 \\
    3x+8y+2z &=19
\end{align*}
which, to allow us to perform Gaussian elimination, can be written in the form:
\begin{align*}
\left[
\begin{array}{rrr|r}
1 & 2 & -1 & 3 \\
2 & 5 & 1 & 12 \\
3 & 8 & 2 & 19
\end{array}
\right],
\end{align*}
where $A = \begin{bmatrix} 1 & 2 & -1 \\
2 & 5 & 1 \\
3 & 8 & 2\end{bmatrix}$.

\textcolor{blue}{Exercise: 
\begin{enumerate}
    \item Determine the row operations necessary to put $A$ into an upper triangular form.  Write these row operations as elementary matrices multiplying $A$ (there should be 3).   
    \item Why are elementary matrices invertible?
    \item Compute $E_1A$, then $E_2E_1A$, then $E_3E_2E_1A$.  What do you notice about the structure of each of the elementary matrices, and the structure of each of the products?
    \item Write the resulting upper triangular matrix as the product $L_2L_1A$, where $L_1$ takes $A$ and zeroes out the first column except for the first entry, and $L_2$ takes $L_1A$ and zeros out the second column except for the first two entries.  (These are called elimination matrices.)
    \item Do everything again, except instead of using numbers, consider a general matrix 
    \begin{align*}
        A = \begin{bmatrix} a_{11} & a_{12} & a_{13} \\
        a_{21} & a_{22} & a_{23} \\
        a_{31} & a_{32} & a_{33}\end{bmatrix}
    \end{align*}
    and assume $a_{ij} \neq 0$ for all $i,j \in \{1,2,3\}$.
\end{enumerate}
}
Answer: \begin{align*}
E_1 &= \begin{bmatrix}
1 & 0 & 0\\
-2 & 1 & 0\\
0 & 0 & 1
\end{bmatrix}\quad
E_2 = \begin{bmatrix}
1 & 0 & 0\\
0 & 1 & 0\\
-3 & 0 & 1
\end{bmatrix}\quad
E_3 = \begin{bmatrix}
1 & 0 & 0\\
0 & 1 & 0\\
0 & -2 & 1
\end{bmatrix}
\end{align*}
\[L_1 = E_2E_1, L_2 = E_3\]


Taxonomy of methods for factoring a matrix: 
Gram-Schmidt: $A = QR$ by triangular orthogonalization
\\
Householder: $A = QR$ by orthogonal triangularization
\\
Gaussian elimination: $A = LU$ by triangular triangularization

\vspace{0.2in}

General formulas for a square matrix:
Let $x_k$ denote the column of $E_{k-1} \cdots E_1A$.  Then $L_k$ must be chosen such that
\begin{align*}
    x_k = \begin{bmatrix}
        x_{1k} \\ \vdots \\ x_{kk} \\ x_{k+1,k} \\ \vdots\\ x_{mk}
    \end{bmatrix} \;
    \stackrel{L_k}{\rightarrow} \; L_kx_k = \begin{bmatrix}
        x_{1k} \\ \vdots \\ x_{kk} \\ 0 \\ \vdots \\0
    \end{bmatrix}
\end{align*}

To do this, we wish to subtract $l_{jk}$ times row $k$ from row $j$, where $l_{jk}$ is the multiplier 
\begin{align*}
    l_{jk} = \frac{x_{jk}}{x_{kk}}, \quad k < j \leq m.
\end{align*}
The matrix $L_k$ then takes the form
\begin{align*}
    L_k = \begin{bmatrix}
        1 & & & & & \\
         & \ddots & & & & \\
         & & 1 &&&\\
         && -l_{k+1,k} & 1 & & \\
         &&\vdots&& \ddots &\\
         && -l_{mk} &&& 1
    \end{bmatrix}
\end{align*}
with the nonzero subdiagonal entries situated in column $k$.

In practice, these $l_{jk}$ multipliers are computed and stored directly, so we neither have to form elementary matrices nor elimination matrices outright.  (See Algorithm 20.1 in your textbook.)

Only problem: this method is not backward stable.  Will fix, next time.

\end{document}



